\documentclass{article}
\usepackage[margin=2cm, headheight=0pt, headsep=1cm, includeheadfoot, top=0.75cm, bottom=1cm]{geometry}
\usepackage{enumerate, fancyhdr, graphicx, amsmath, float, url, hyperref, color}
\usepackage{array,booktabs}

\title{Vault 5431 - Assurance}
\author{Alicia Wu, Britney Wong, Chang Yang Jiao, Paul Chesnais}
\date{\today}

\pagestyle{fancy}
\fancyhead{}
\lhead{cj285, bmw227, pmc85, yw344}
\chead{Vault 5431 - Assurance}
\rhead{\today}
\fancyfoot{}
\rfoot{\thepage}
\lfoot{\includegraphics[height=20pt]{figures/Logo}}
\renewcommand{\headrulewidth}{0.5pt}
\renewcommand{\footrulewidth}{0.5pt}

\begin{document}
\maketitle
\thispagestyle{empty}
\section{Alpha}
\par Every method written for this release was tested using JUnit tests. In addition to functionality related to Auditing, all of the boilerplate related to encryption and disk I/O was extensively tested. Due to the concurrent nature of a web server, all methods handling file access were carefully designed to be thread safe. Sadly, this is much harder (if not impossible) to test, and these methods were only tested in a serial manner. But, there is reason to believe that because proper design principles with respect to concurrency were employed when building this and that as long as each file is accessed using the proper methods, race conditions will be avoided.

\par Additionally, the codebase was checked using FindBugs. There are no bugs that are deemed ``Scary'' or ``Scariest''. The build manager checksums all dependencies and verifies before using them, but the Bouncy Castle cryptographic library is pulled in and checked manually, along with the build manager itself at first install. There is a final web dependency that needed to be edited and was compiled directly from source, and is therefore trusted.

\par In the alpha, we also made a user interface that the client can interact with. We tested the forms for odd or malformed inputs, tested every page to ensure that there are no unexpected outputs, and ensured that the server handled incorrect inputs or requests correctly. For example, the server warns the user when inputs are incorrect and fails gracefully when receiving such inputs by responding with the associated error code (i.e. not 500 Internal Error).

\section{Beta}
\par As with the Alpha release, JUnit tests were written and ran for every method. As new classes and methods were created and implemented, old and new JUnit tests were ran to ensure that all code is running as expected. For 2FA with text messages, not only did we test with JUnit using a demo user and demo number, but we also tested extensively using real phone numbers, making sure that real phones could receive the authentication code sent by the server via Twilio.

\par Additionally, the codebase was checked using FindBugs. There are no ``Scariest'' and ``Scary'' Bugs. ``Troubling'' bugs were found periodically and fixed. At the end of the sprint, there were two "Troubling" bugs of very low rank, but were quickly either fixed or deemed not a bug. Additionally, the build manager checksums all dependencies and verifies before using them. The server also had to manually add Twilio's SSL certificate to the trust store but since it is trusted and reputable, this is okay.

\par In the Beta, we added onto the Alpha user web interface. We tested every form for bad input and did form validation on both the client side and the server side. We also tested the forms and routes for input that is entered via Ajax POST or GET calls. The server side validation will catch errors made through these calls. We also tested UI functionality written in Javascript by extensively interacting with the web interface via clicks and refreshes.

\end{document}
